% Source fingerprint: ff5e9a5fcb885efbd83770ad0a51e844e6755106c78767056860332571eeb2b1
% T06: raw TeX cells preserved; no numeric highlighting.
\begin{table}[htbp]
\centering\small\renewcommand{\arraystretch}{1.18}\setlength{\tabcolsep}{4pt}
\caption[\(L^p\) 稠密逼近工具选择]{\(L^p\) 稠密逼近工具选择。无穷端点仍可用简单函数按本性上确界逼近，但连续、紧支集或光滑稠密性不能照搬。}
\label{b1:tab:visual-t06}
\begin{tabularx}{\linewidth}{@{}>{\raggedright\arraybackslash}p{3.1cm}>{\raggedright\arraybackslash}p{4.7cm}>{\raggedright\arraybackslash}X@{}}
\toprule
逼近对象 & 适用范围 & 必须保留的条件 \\
\midrule
简单函数 & 任意测度空间；\(1\le p\le\infty\) & 简单是有限值域，不等于非零区域测度有限 \\
有限测度非零区域的简单函数 & 任意测度空间；\(1\le p<\infty\) & 有限测度区域由待逼近函数本身选择，不要求底测度先为有限 \\
有界截断和空间截断 & 有限 \(p\)；按正文的区域穷竭使用 & 分别控制幅值尾部与空间尾部；无穷测度空间上两者不能混同 \\
\(C_c(\mathbb R^d)\) & Lebesgue 测度；\(1\le p<\infty\) & 先用正则性逼近可测集，再构造连续截断 \\
\(C_c^\infty(\mathbb R^d)\) & Lebesgue 测度；\(1\le p<\infty\) & 继续使用第20章的磨光；不能在一般测度空间凭空要求光滑 \\
\bottomrule
\end{tabularx}
\end{table}
